\section{The arithmetic bound and the Riemann hypothesis}\label{sec:arithmetic-bound}

It suffices to bound the excess at points spaced $O(x^{3/4})$ apart.
The convexity of $G(y)+y\log y$ controls the error between them.

At a point $x$ put $a=xb$ and keep $Q=N_Tmb\delta_2/2$ from
\S\ref{sec:weights}. Let $m$ be the number of lifts and $k$ the
positive index of the first hermitian space. The rank choices in
\citep[Proposition~5.15]{adriazola-grh} satisfy
\begin{equation}\label{eq:arithmetic-ranks}
  x^{1/2-\varepsilon}\le m\le x^{1/2+\varepsilon/2},
  \qquad xm\le k\le n^2x^{1/2-\varepsilon}.
\end{equation}
Write
\[
  \mathcal B=N_Ta\left(1+\frac{m^2}{k}\right)
       +m\left(2nmbx^\varepsilon+\frac{N_2kb x^{2\varepsilon}}{n^2}\right).
\]

\begin{hypothesis}[\openinput]\label{hyp:one-shift-arithmetic}
For arbitrarily small $0<\varepsilon<1/52$ there are constants
$J\ge3/2$ and $K>0$ with the following property for arbitrarily
large $Y$. There are primes $q,d$ as in
Proposition~\ref{prop:field-window} on a window $[U,U^J]$ containing
$[Y,Y^{3/2}]$, and there is a finite mesh
\[
  Y=x_0<\cdots<x_\nu=Y^{3/2},\qquad
  x_{i+1}-x_i\le Kx_i^{3/4}\quad(0\le i<\nu).
\]
At each mesh point $x$ there are ranks satisfying
\eqref{eq:arithmetic-ranks}, a set $T_2$ with $N_2\ge nx^{-\varepsilon}$,
and a weight set $T$ satisfying the size and translate bounds of
Lemma~\ref{lem:weight-count}. For this $T$ there are a shift $u_0\in H$,
weights $(a_t)_{t\in T}$ and nodes satisfying \eqref{eq:weight-ranges}
such that
\begin{equation}\label{eq:arithmetic-bound}
  D_{u_0}\le x^{5\varepsilon}\mathcal B .
\end{equation}
\end{hypothesis}

The primes stay fixed on a window. The sets, ranks, and shift may
depend on the mesh point. To prove \eqref{eq:arithmetic-bound} by
Proposition~\ref{prop:accounting}, it suffices to realize
\eqref{eq:acct-gain}--\eqref{eq:acct-base} on one chain with these
weights and nodes and to bound its total error by
\[
  e_1+e_2+e_3+e_4+e_5+A e_6+e_7\le x^{5\varepsilon}\mathcal B,
  \qquad A=\sum_{t\in T}a_t .
\]

\begin{theorem}\label{thm:arithmetic-rh}
Hypothesis~\ref{hyp:one-shift-arithmetic} implies the Riemann hypothesis.
\end{theorem}

\begin{proof}
Fix an $\varepsilon$ from the hypothesis and write
$\theta=1/2+8\varepsilon<1$. At a mesh point, division by $Q$ gives
\[
  \frac{\mathcal B}{Q}
   =\frac{2x}{\delta_2}\left(\frac1m+\frac m k\right)
     +\frac{4nm x^\varepsilon}{N_T\delta_2}
     +\frac{2k x^{2\varepsilon}}{nN_T}.
\]
Since $\delta_2\ge x^{-\varepsilon}$ and $n/N_T\le2x^{\varepsilon/8}$,
\eqref{eq:arithmetic-ranks} bounds the right side by
\[
  2x^{1/2+2\varepsilon}+2x^\varepsilon
  +8x^{1/2+21\varepsilon/8}+4x^{1/2+9\varepsilon/8}
  \ll x^{1/2+3\varepsilon}.
\]
By \eqref{eq:one-shift}, \eqref{eq:field-projection} and
\eqref{eq:field-window-bound}, the chosen shift therefore gives
\[
  G(x)\le\frac{D_{u_0}}{Q}+3x^{1/2+\varepsilon}+o(1)
       \ll_{\varepsilon} x^\theta .
\]
Here \eqref{eq:arithmetic-bound} bounds the first term, and
Proposition~\ref{prop:one-shift} bounds the error uniformly on the mesh.

Take consecutive mesh points $u<v$ and let $L$ denote affine
interpolation at $u,v$. Put $f(y)=y\log y$. By \eqref{eq:prime-sum},
\[
  G(y)+f(y)=(1+\gamma)y+\sum_{r\ge1}\Lambda(r)(y/r-1)_+
\]
is convex. Thus $G(y)\le LG(y)+Lf(y)-f(y)$ on $[u,v]$.
The mesh estimate and concavity of $y^\theta$ give
$LG(y)\le C_\varepsilon y^\theta$. Since $f''(y)=1/y\le1/u$,
the function $f(y)-y^2/(2u)$ is concave, whence
\[
  Lf(y)-f(y)\le\frac{(y-u)(v-y)}{2u}
        \le\frac{(v-u)^2}{8u}\le\frac{K^2}{8}\sqrt y .
\]
Consequently $G(y)\le(C_\varepsilon+K^2/8)y^\theta$ throughout
$[Y,Y^{3/2}]$, with one constant for all the chosen windows.

Theorem~\ref{thm:window-criterion} now excludes zeros of $\zeta$ in
$\operatorname{Re}s>1/2+8\varepsilon$. The hypothesis admits
arbitrarily small $\varepsilon$, so every nontrivial zero $\rho$
satisfies $\operatorname{Re}\rho\le1/2$. The functional equation gives
the same inequality for $1-\rho$, hence $\operatorname{Re}\rho=1/2$.
\end{proof}
